\documentclass[a4paper,11pt]{article}

\usepackage[utf8]{inputenc}
\usepackage[czech]{babel}
\usepackage[left=2cm,top=3cm,text={17cm,24cm}]{geometry}
\usepackage{graphicx}
\usepackage{listings}
\usepackage{url}

\title{ISA - Síťové aplikace a správa sítí\\
{\bf\large Laboratorní manuál}}

\author{Vysoké učení technické v Brně}

\date{\url{https://github.com/nesfit/ISA/tree/master/manual}}

\begin{document}

{\let\newpage\relax\maketitle}

Tento laboratorní manuál slouží jako referenční příručka pro laboratorní cvičení
předmětu ISA -- Síťové aplikace a správa sítí. Očekává se, že si studenti před
účastí na jednotlivých cvičeních přečtou sekce potřebné pro dané cvičení.

Sekce \ref{adresy_ipv4}, \ref{adresy_ipv6}, \ref{basics}, \ref{wireshark} a \ref{ipconfig} jsou
potřebné pro všechny cvičení. Cvičení \emph{Zabezpečený přenos dat} navíc
předpokládá znalosti ze sekce~\ref{ntp}, \ref{ssh} a~\ref{tls}. Ve cvičení \emph{DNS,
DNSSEC} se využívají sekce \ref{dns} a \ref{dnssec}. Cvičení
\emph{Konfigurace a analýza přenosů VoIP} předpokládá znalosti uvedené v
sekci~\ref{sip}. Cvičení \emph{Správa a monitorování sítě} staví na informacích
uvedených v sekcích \ref{syslog} a \ref{netflow}.

\setcounter{tocdepth}{1}
\tableofcontents
\newpage

\section{Adresy v~IPv4 sítí}\label{adresy_ipv4}

\subsection{Formát adresy}
IPv4 adresa je délky 32 bitů a preferovaný zápis má formát {\tt X.X.X.X}, kde
{\tt X} je decimální zápis 8 bitového čísla. Příklad:\\
\begin{verbatim}
8.8.8.8
127.0.0.1
\end{verbatim}

Adresa dvě části:
\begin{itemize}
    \item adresa sítě ({\bf prefix}) a
    \item adresa uzlu.
\end{itemize}
Délka prefixu se zapisuje v~desítkovém tvaru za lomítko. {\tt 192.168.0.1/24}

V~síti s~daným prefixem existují 2 speciální adresy, které nelze použít pro
adresování jednotlivých uzlů:
\begin{itemize}
    \item adresa sítě - adresa uzlu je nulová, např. 192.168.0.0/24
    \item broadcastová adresa - nejvyšší možná adresa uzlu pro daný prefix,
        např. 192.168.0.255/24
\end{itemize}

\section{Adresy v~IPv6 sítí}\label{adresy_ipv6}

\subsection{Formát adresy}
IPv6 adresa je délky 128 bitů a zapsaných ve formátu {\tt X:X:X:X:X:X:X:X}, kde
{\tt X} je hexadecimální zápis 16 bitového čísla.
Příklad:\\
\begin{verbatim}
FEDC:BA98:7654:3210:fedc:ba98:7654:3210
1080:0000:0000:0000:0008:0800:200C:417a
\end{verbatim}

Preferovaný formát je dále upraven v~RFC 5952\footnote{https://tools.ietf.org/html/rfc5952}, které definuje následující
pravidla:
\begin{itemize}
    \item Nuly na začátku každého 16 bitového čísla je potřeba vynechat (např.
        80 namísto 0080 a 0 namísto 0000).
    \item Více nulových bloků lze nejvýše jednou nahradit znakem {\tt ::}, a
        MUSÍ být nahrazena nejdelší možná posloupnost takových nulových bloků.
        V~případě více shodně dlouhých posloupností, nahrazuje se ta nejvíce
        vlevo.
    \item Znak {\tt ::} nesmí nahrazovat samostatný nulový blok.
    \item Znaky "a", "b", "c", "d", "e", "f" hexadecimální soustavy se vždy
        píšou malými písmeny.
\end{itemize}
Adresy z~předchozího příkladu tedy budou vypadat následovně:\\
\begin{verbatim}
fedc:ba98:7654:3210:fedc:ba98:7654:3210
1080::8:800:200c:417a
\end{verbatim}

Stejně jako v~IPv4 má adresa dvě části: \emph{adresa sítě ({\bf prefix})} a
\emph{adresa uzlu}. Délká prefixu se zapisuje v~desítkovém tvaru za lomítko
(např. {\tt 1080::/60}).

\subsection{Rozdělení adres}
IPv6 adresy je možno rozdělit podle rozsahu. Typicky může jít o~tři
možnosti -- adresy na lince (neprojdou za router), lokální adresy (ULA, nejsou
routovatelné ve veřejné sítí) a veřejné adresy.

\begin{table}[ht!]
    \begin{center}
        \begin{tabular}{l|l|l}
            Význam & Prefix (bitově) & Prefix \\
            \hline\hline
            Veřejné  & {\tt 001} & {\tt 2000::/3} \\
            \hline
            Lokální  & {\tt 1111 110} & {\tt FC00::/7} \\
            \hline
            Linkové  & {\tt 1111 1110 10} & {\tt FE80::/10} \\
            \hline
            Multicast & {\tt 1111 1111} & {\tt FF00::/8} \\
            \hline
        \end{tabular}
    \end{center}
\end{table}

\subsection{Lokální IPv6 (ULA) adresy}\label{ula}
Síťová část ULA adresy se stává ze 4 části:
\begin{description}
    \item [Prefix] {\tt fc00::/7}
    \item [L] bit 1 pokud byl prefix přiřazen lokálně, 0 zatím nebyla
        definována, začátek adresy je proto typicky {\tt fd00::/8}
    \item [Global ID] identifikátor sítě, měl by být unikátní, standard
        popisuje pseudonáhodný algoritmus pro generování. Unikátnost je
        požadována, aby při spojení více lokálních sítí nebylo třeba žádnou
        přečíslovat.
    \item [Inderface ID] Identifikátor rozhraní, existuje několik variant jak
        jej získat. Původní standard (RFC
        3513\footnote{https://tools.ietf.org/html/rfc3513} a RFC
        4291\footnote{https://tools.ietf.org/html/rfc4291}) doporučoval použití
        EUI-64. V~současnosti byla tato metoda nahrazena RFC
        7217\footnote{https://tools.ietf.org/html/rfc7217}. Pro ochranu
        soukromí jsou další specifika automatické generace identifikátorů
        definovány v~RFC 4941\footnote{https://tools.ietf.org/html/rfc4941}.
\end{description}

\begin{table}[ht!]
    \begin{center}
        \begin{tabular}{c|c|c|c|c}
            7 bits & 1 &  40 bits  &  16 bits  & 64 bits \\
            \hline
            1111 110 & L & Global ID & Subnet ID & Interface ID \\
            \hline
        \end{tabular}
    \end{center}
\end{table}

Lokální síť je tedy složená ze tří části:
\begin{itemize}
    \item unikátní prefix délky 48 bitů,
    \item 16 bitový identifikátor podsítě,
    \item adresa uzlu o~délce 64 bitů.
\end{itemize}

\section{Základy konfigurace linuxového serveru} \label{basics}
Při práci s unixovými systémy doporučujeme využívat manuálové stránky pro zjištění významu daného příkazu (aplikace) a detailního popis parametrů spuštění, viz příkaz {\tt man <příkaz>}, např. {\tt man ifconfig}.
\subsection{Základní orientace v~Linuxu}
Většina práce v~OS Linux bude probíhat v terminálovém okně. Terminál na školních PC spustíte pomocí \texttt{Alt+F2}, zde zadáte "\texttt{gnome-terminal}", případně můžete vybrat aplikaci terminálu v nabídce aplikací.

\subsubsection{Unixový souborový systém}
Všechny soubory v počítači jsou uloženy v~souborovém systému. Ten je organizován jako invertovaný strom souborů a adresářů, kde kořenový adresář \texttt{root} je označen jako lomítko  \texttt{"/"}. Kořenový adresář obvykle obsahuje standardizovaný seznam podadresářů, do kterých se ukládají systémové, aplikační i uživatelské soubory. 

Některé podadresáře a jejich význam:
\begin{itemize}
  \item {\tt /etc} - konfigurační soubory operačního systému,
  \item {\tt /bin} - systémové aplikace a příkazy
  \item {\tt /sbin} - systémové aplikace a příkazy pro správce systému (uživatel {\tt root})
  \item {\tt /usr} - uživatelské aplikace a příkazy ({\tt /usr/bin, /usr/lib, /usr/local})
  \item {\tt /home} - domovské adresáře uživatelů
  \item {\tt /var} - soubory s proměnlivým obsahem, cache paměti, logy (\texttt{/var/log}),
  \item {\tt /dev} - ovladače (drivers) k periferním zařízením počítače.
\end{itemize}

\subsubsection{Základní příkazy pro práci se soubory, adresáři a procesy}
Některé základní příkazy pro práci v~terminálu OS Linux:
\begin{itemize}
  \item \texttt{cd} - změna adresáře,
  \item \texttt{ls} - zobrazení obsahu adresáře,
  \item \texttt{cp} - kopírování souborů,
  \item \texttt{mv} - přesun souborů,
  \item \texttt{mkdir} - vytvoření adresáře,
  \item \texttt{rmdir} - zrušení prázdného adresáře,
  \item \texttt{touch} - vytvoření prázdného souboru,
  \item \texttt{rm} - smazání souboru,
  \item \texttt{head} - zobrazení začátku souboru,
  \item \texttt{tail} - zobrazení konce souboru,
  \item \texttt{cat} - zobrazení obsahu souboru,
  \item \texttt{more} - výpis obsahu souboru po stránkách,
  \item \texttt{grep} - hledání textu v souboru,
  \item \texttt{chmod} - změna přístupových práv souboru,
%  \item \texttt{sed} - editace textu v~příkazové řádce,
%  \item \texttt{awk} - skenování a zpracování textu,
  \item \texttt{ps} - výpis běžících procesů.
\end{itemize}

Mezi textové editory patří například \texttt{nano}, {\tt vi} nebo \texttt{vim}. Tyto editory je vhodné používat k editaci konfiguračních soborů v operačním systému. Kromě řádkových textových editorů lze použít použít i grafické editory, například \texttt{gedit}.

\subsubsection{Role uživatele a správce v OS}
Každý uživatel operačním systému Linux má svůj jedineční identifikátor \texttt{uid} (user ID). Kromě uživatelů jsou v OS definovány i skupiny uživatelů definované identifikátor {\tt gid} (group ID). Každý soubor je vlastněn konkrétním uživatelem patřící do určité skupinou. Přístup k souboru je omezen přístupovými právy pro vlastníka souboru, skupinu či ostatní uživatele. 

Každý proces (běžící aplikace) v~systému je spuštěn pod určitým uživatelem a je jednoznačně identifikován číslem procesu {\tt pid} (process ID).

Základní role uživatelů v unixovém OS:
\begin{itemize}
  \item \texttt{root} - správce systému: má veškerá práva a kontrolu nad systémem.
  \item \texttt{user} - běžný uživatel: má pouze základní správa bez možnosti zasahovat do systémového nastavení.
%  \item \texttt{sudo} - delegace některých práv super uživatele na normálního uživatele (nastavení v~\texttt{/etc/sudoers}).
\end{itemize}

Při práci v operačním systému se snažíme vyhnout používání úrovně  \texttt{root}. Pokud potřebujeme použít pro určité činnosti příkazy vyžadující úroveň práv {\tt root}, můžeme využít příkaz \texttt{sudo}, který deleguje práva správce systému danému uživateli bez nutnosti přihlásit se jako uživatel {\tt root}. Toto oprávnění je definováno v souboru {\tt /etc/sudoers}. Příkaz s privilegovaným přístupem pak můžeme spustit příkazem {\tt sudo command}. 

V případě potřeby je možné změnit roli uživatele příkazem \texttt{su [user]}
\begin{itemize}
  \item \texttt{su} - přepnutí na uživatele \texttt{root}.
  \item \texttt{su user} - přepnutí na uživatele \texttt{user}.
\end{itemize}

\subsection{Správa systémových služeb}
\label{sluzby}
Operační systém Linux poskytuje řadu nástrojů pro správu systémových služeb. 
%různých aplikací plnících roli systémových služeb, tyto mohou poskytovat užitečné funkce jiným aplikacím, uživatelům nebo spravovat konfiguraci subsystémů.
Systémové služby obvykle běží autonomně, na pozadí systému bez přímého rozhraní pro uživatele. Operační systém proto poskytuje nástroje pro jejich manipulaci. Správu operačního systému a systémových služeb v linuxových distribucích provádí aplikace \texttt{systemd} (system demon). Pro spouštění a manipulaci systémových služeb se využívá nástroj {\tt systemctl} (system control).

Příklady použití: 
\begin{itemize}
  \item \texttt{systemctl start [služba]} -- spustí službu
  \item \texttt{systemctl stop [služba]} -- zastaví službu
  \item \texttt{systemctl restart [služba]} -- restartuje službu
  \item \texttt{systemctl enable [služba]} -- povolí automatické spuštění služby po startu systému
  \item \texttt{systemctl disable [služba]} -- zakáže automatické spuštění služby po startu systému
  \item \texttt{systemctl status [služba]} -- vypíše informace aktuálního stavu služby a několik posledních řádků systémových logů této služby. Pokud spustíme příkaz bez parametrů, vypíš  informace
        o~všech aktuálně spuštěných službách.
\end{itemize}

Jméno služby má typicky tvar \texttt{<aplikace>.service}, například
\texttt{sshd.service}. Pro zobrazení logů konkrétní systémové služby lze použít příkaz \texttt{journalctl -u [služba]}. Detailní popis příkazů a jejich parametrů naleznete v~manuálových stránkách.

\subsection{Konfigurace síťových zařízení} \label{basic_ipconfig}
V této kapitole si ukážeme několik nástrojů operačního systému Linux, které slouží pro zjišťování síťové konfigurace.

\subsubsection{Zobrazení konfigurace}
Základním příkazem je příkaz \texttt{ip}. Pomocí tohoto příkazu můžeme zjistit konfiguraci všech síťových rozhraní počítače (fyzických i virtuálních), obsah směrovací tabulky a jiné. Manuálové stránky pro dané možnosti zobrazíte jako \texttt{man ip <volba>}, např. {\tt man ip link}.

\begin{itemize}
  \item \texttt{ip address} -- zobrazí adresu na všech síťových rozhraní (též příkaz {\tt ifconfig}). 
  \item \texttt{ip route} -- zobrazí pravidla ve směrovací tabulce
  \item \texttt{ip link} -- zobrazí konfiguraci síťových rozhraní
  \item \texttt{ip neighbour} -- zobrazí obsah ARP záznamů (též příkaz \texttt{arp}).
\end{itemize}

Příkaz {\tt ip} neslouží jen k výpisu konfiguraci, ale je možné jím i nastavit parametry síťového rozhraní.

Další příkazy pro zobrazení síťové konfigurace:
\begin{itemize}
  \item \texttt{netstat -rn} -- zobrazí obsah směrovací tabulky
  \item \texttt{ss} (\texttt{sockstat}) -- zobrazí čekajících spojení či navázaných síťových spojení včetně zobrazení typu transportního protokolu (parametry {\tt -t} pro TCP, {\tt -u} pro UDP), zdrojových a cílových adres a portů.
  \item {\tt ss state ESTABLISHED -t} -- zobrazí navázaná spojení nad TCP (ve stavu ESTABLISHED)
\end{itemize}

\subsubsection{Testování konektivity a síťových služeb}
Následující příkazy slouží k testování dostupnosti zařízení a konektivity. 
\begin{itemize}
  \item \texttt{ping <ipv4> | <hostname>} -- zašle příkaz ICMP ECHO\_REQUEST na dané zařízení.
  \item \texttt{ping6 <ipv6> | <hostname>} -- zašle příkaz ICMPv6 ECHO\_REQUEST na dané zařízení.
  \item \texttt{traceroute <ipv4> | <hostname>} -- otestuje a zobrazí cestu IP datagramu  k cíli.
  \item \texttt{traceroute6 <ipv6> | <hostname>} -- otestuje a zobrazí cestu IPv6 datagramu k cíli.
  \item \texttt{tcptraceroute <ipv4> | <ipv6> | <hostname>} - otestuje a zobrazí cestu paketu TCP k cíli.
\end{itemize}

\noindent
Pro přihlášení na vzdálený host můžeme použít dva nástroje:
\begin{itemize}
  \item \texttt{telnet <host>} -- vzdálený přístup pomocí nešifrovaného spojení
  \item \texttt{ssh <host>} -- vzdálený přístup pomocí šifrovaného přípojení
\end{itemize}

\noindent
Testování dostupných zařízení a služeb
\begin{itemize}
  \item \texttt{nmap} - skenování aktivních stanic v síti, otevřených portů apod. 
  \item \texttt{telnet <host> <port>} -- testování dostupnosti služeb nad TCP
\end{itemize}

\subsection{Konfigurační soubory}
Konfigurační soubory operačního systému a služeb se většinou ukládají do adresáře \texttt{/etc}. Jednotlivé služby zde mají svůj konfigurační soubor s názvem dané služby. Konfigurace některých služeb je ulžena v adresářích  \texttt{"/etc/<sluzba>.d/"} (například \texttt{/etc/rsyslog.d/}). 

Důležité  konfigurační soubory:
\begin{itemize}
  \item \texttt{/etc/hostname} -- obsahuje doménové jméno linuxového systému.
  \item \texttt{/etc/hosts} -- obsahuje statické mapování doménového jména na IP adresu (nevyžaduje DNS)
  \item \texttt{/etc/host.conf} -- obsahuje konfiguraci specifickou pro DNS resolver.
  \item \texttt{/etc/resolv.conf} -- konfigurační soubor pro DNS resolver, který obsahuje zejména odkazy na lokální DNS servery (získané z DHCP)
  \item \texttt{/etc/rsyslog.conf} - konfigurační soubor aplikace {\tt rsyslog} 
\end{itemize}
Podrobnosti o konfiguraci výše uvedených služeb lze získat z manuálových stránek, např. {\tt man resolv.conf}. 

\subsection{Logování událostí}
Systémové i aplikační události jsou zaznamenávány do logovacích souborů. Logovací soubory jsou obvykle uloženy v~adresáři \texttt{/var/log}. Tyto soubory jsou cyklicky rotovány po určitém čase či po dosažení určité velikosti. Dobu, po kterou jsou logy udržovány, lze nastavit v~konfiguračním souboru\\\texttt{/etc/logrotate.conf}.

Důležité logovací soubory:
\begin{itemize}
  \item \texttt{/var/log/messages} -- obsahuje globální systémové zprávy a události.
  \item \texttt{/var/log/auth.log} -- obsahuje zprávy týkající se autentizace uživatelů.
  \item {\tt ./var/log/maillog} -- obsahuje logovací zprávy od mailového serveru. 
%  \item \texttt{/var/log/kern.log} -- ukládá zprávy týkající se jádra OS.
\end{itemize}

Tyto logovací soubory můžeme prohlížet přímo pomocí standardních zobrazovacích příkazů ({\tt cat, more, tail -f, grep}). Zobrazení logů většinou vyžaduje přístupová práva správce systému. 

Systémové události můžeme také prohlížet pomocí nástroje \texttt{journalctl}, který umožňuje podrobněji procházet seznam zpráv, filtrovat a vyhledávat.

Příklad použití \texttt{journalctl}:
\begin{itemize}
  \item \texttt{journalctl -n <x>} - zobrazí posledních x záznamů.
  \item \texttt{journalctl -p <priority>} - zobrazí pouze záznamy s~danou prioritou Syslog.
  \item \texttt{journalctl -u <unit> | <pattern>} - zobrazí pouze záznamy s danou službou nebo obsahující zadaný řetězec
  \item \texttt{journalctl -t <syslog\_id> | <pattern>} -- vyhledává záznamy obsahující zadané ID syslogu nebo záznamy obsahující v ID syslogu zadaný řetězec.
  \item \texttt{journalctl -f} - zobrazí posledních deset událostí.
  \item \texttt{journalctl --since <whenstart> --until <whenend>} - zobrazí záznamy v~daném časovém rozmezí.
\end{itemize}






\section{Síťový analyzátor Wireshark}\label{wireshark}
Wireshark je grafická aplikace sloužící k zachytávání paketů a jejich analýze, viz obrázek \ref{fig:wireshark-layout}. Wireshark umí zachytávat provoz na lokálních fyzických či virtuálních síťových rozhraních počítače nebo zpracovávat zachycené pakety uložené v souboru typu PCAP. V případě zpracování příliš velkých souborů (stovky MB data) je vhodnější využívat pro zpracování řádkovou verzi nástroje Wireshark, která se nazývá {\tt tshark}\footnote{Viz \url{https://www.wireshark.org/docs/man-pages/tshark.html}}. Další populární nástroj pro odchytávání paketů a jejich zpracování je unixová aplikace \texttt{tcpdump}. 

\begin{figure}[h]
  \centering
  \includegraphics[width=170mm]{fig/wireshark-layout.pdf}
  \caption{Wireshark}\label{fig:wireshark-layout}
\end{figure}

\subsection{Zachytávání provozu (Packet Capturing)}
Následující kroky popisují proces zachytávání paketů v programu Wireshark:
\begin{enumerate}
  \item Zachytávat provoz můžeme z vybraného síťového rozhraní nebo z~více rozhraní najednou. Rozhraní vybereme v menu \texttt{Capture -> Options}.
  \item V nastavení Wiresharku je možné nastavit filtr pro výběr zachytávaných paketů (Capture Filter), pomocí kterého řekneme aplikaci, jaké pakety má zachytávat. Protože zpracování každého paketu je výpočetně náročné (dochází k parserování jednotlivých vrstev paketu a detekci protokolů), snížíme použití vstupního filtru výpočetní nároky na odchyt paketů. Políčko {\tt Capture Filter} umožňuje vybrat již předdefinovaný filtr nebo je zadat vlastní filtr, viz {\tt man pcap-filter}.
  \item Záchyt spustíme pomocí tlačítka \texttt{Start}.
  \item Odchycený provoz (tj. jednotlivé pakety spolu s časovými značkami) můžeme také uložit do souboru typu PCAP pro pozdější analýzu.
\end{enumerate}

\subsection{Výběr zachyceného provozu pomocí zobrazovacího filtru (Display Filter)}
Zachycené a zpracované pakety jsou zobrazeny v okně {\tt Packet List Pane}. Protože nás často zajímá pouze určitý provoz, např. komunikace HTTP s konkrétní stanicí, je vhodné využít zobrazovací filtr (Display Filter), který se nachází pod nástrojovou lištou. Pomocí tohoto filtru vybereme pouze pakety s požadovanou vlastností, např. filtr {\tt ip.addr == 147.229.8.12} vybere veškeré pakety, které mají danou zdrojovou či cílovou IP adresu. 

Základní operace pro vytváření zobrazovacího filtru je popsána v tabulce \ref{tab:display-filter}. 
\begin{center}
  \begin{table}[h]
    \centering
    \def\arraystretch{1.2}
    \begin{tabular}{|l|l|}
      \hline
      \textbf{Porovnávání} & \texttt{==, >=, <=, !=, contains}\\
      \textbf{Logické operátory} & \texttt{||, or, \&\&, and, !, not}\\
      \textbf{Kombinace filtrů} & \texttt{(ip.src==192.168.0.105 and udp.port==53)}\\
      \textbf{Filtrování na základě existence pole} & \texttt{http.cookie or http.set\_cookie}\\
      \textbf{Filtrování specifických bytů} & \texttt{eth.src[4:2]==22:1b}\\
      \textbf{Použití regulárních výrazů} & \texttt{http.host \&\& !http.host matches ".com"}\\
      \hline
    \end{tabular}
    \caption{Syntax zobrazovacího filtru (Display Filter)}\label{tab:display-filter}
  \end{table}
\end{center}

\subsection{Zobrazení toků (Flow Graph)}
Wireshark umožňuje zobrazit časovou posloupnost komunikace pomocí tzv. grafu toků zvaného \texttt{flow graph}, viz menu \texttt{Statistics -> Flow Graph}. Graf toků zobrazuje komunikaci mezi jednotlivými zařízeními definovanými IP adresou, viz obrázek \ref{fig:flow-graph}. Pomocí grafu toků můžeme vidět, jak často spolu stanice komunikují, jaké zprávy a v jakém pořadí si posílají a další.

\begin{figure}[h]
  \centering
  \includegraphics[width=110mm]{fig/flow-graph.png}
  \caption{Flow Graph}\label{fig:flow-graph}
\end{figure}

\subsection{Zobrazení obsahu toku (Stream Analysis)}
Wireshark zachytává jednotlivé pakety. Pokud chceme sledovat průběh celé komunikace včetně obsahu, můžeme  zobrazit danou komunikaci (streamu) pomocí volby \texttt{Flow <protokol> Stream}, např. TCP či UDP spojení, HTTP spojení apod. Zobrazení toku provedeme tak, že vybereme jeden zachycený paket daného toku, klikneme na něj pravým tlačítkem myši a vybereme volbu  \texttt{Follow <protokol> Stream}. Příklad obsahu HTTP komunikace je na obrázku \ref{fig:follow_stream}.

\begin{figure}[h]
  \centering
  \includegraphics[width=110mm]{fig/follow_stream.png}
  \caption{Zobrazení toku HTTP}\label{fig:follow_stream}
\end{figure}

\subsection{Časové značky}
Každý paket obsahuje časovou značku svého záchytu. Zobrazení časové značky závisí na nastavení Wiresharku. Wireshark umožňuje zobrazit čas jako časovou značku obsahující datum a čas záchytu, absolutní čas, relativní čas od začátku záchytu komunikace a další. Formát času můžeme nastavit v menu \texttt{View -> Time Display Format}, viz obrázek \ref{fig:time-stamps}. Pro některé účely je vhodnější vidět pouze relativní čas od začátku záchytu paketů, pro jiné je potřeba znát absolutní čas záchytu. 

\begin{figure}[h]
  \centering
  \includegraphics[width=110mm]{fig/time-stamps.png}
  \caption{Nastavení formátu časových značek}\label{fig:time-stamps}
\end{figure}

\setcounter{section}{0}
\section*{Teorie}
\section{Správa systémových služeb}\label{sluzby-teorie}
OS Linux poskytuje řadu různých aplikací plnících roli systémových služeb, tyto
mohou poskytovat užitečné funkce jiným aplikacím, uživatelům nebo spravovat
konfiguraci subsystémů. Systémové služby typicky běží autonomně, na pozadí
systému bez přímého rozhraní pro uživatele. Naopak systém poskytuje speciální
rozhraní pro jejich manipulaci. Pro Linuxové distribuce se \texttt{systemd} máme
dostupnou aplikaci \texttt{systemctl}, která poskytuje několik základních
příkazů pro manipulaci systémových služeb:
\begin{itemize}
    \item \texttt{systemctl start [služba]} spustí službu
    \item \texttt{systemctl stop [služba]} zastaví službu
    \item \texttt{systemctl restart [služba]} restartuje službu
    \item \texttt{systemctl enable [služba]} povolí automatické spuštění služby
        po startu systému
    \item \texttt{systemctl disable [služba]} zakáže automatické spuštění
        služby po startu systému
    \item \texttt{systemctl status [služba]} vypíše informace aktuálního stavu
        služby a několik posledních řádků systémových logů této služby. Tento
        příkaz můžete použít i bez specifikace služby, dostanete tak informace
        o~všech aktuálně spuštěných službách.
\end{itemize}

Jméno služby má typicky tvar \texttt{<aplikace>.service}, například
\texttt{sshd.service}.

Pro zobrazení kompletních logů konkrétní systémové služby můžete použít příkaz
\texttt{journalctl -u [služba]}.

Detailní popis včetně možných voleb a příkladů použití můžete nalézt
v~manuálových stránkách.

\section{Adresy v~IPv4 sítí}\label{ipv4-teorie}

\subsection{Formát adresy}
IPv4 adresa je délky 32 bitů a preferovaný zápis má formát {\tt X.X.X.X}, kde
{\tt X} je decimální zápis 8 bitového čísla. Příklad:\\
\begin{verbatim}
8.8.8.8
127.0.0.1
\end{verbatim}

Adresa dvě části:
\begin{itemize}
    \item adresa sítě ({\bf prefix}) a
    \item adresa uzlu.
\end{itemize}
Délka prefixu se zapisuje v~desítkovém tvaru za lomítko. {\tt 192.168.0.1/24}

V~síti s~daným prefixem existují 2 speciální adresy, které nelze použít pro
adresování jednotlivých uzlů:
\begin{itemize}
    \item adresa sítě - adresa uzlu je nulová, např. 192.168.0.0/24
    \item broadcastová adresa - nejvyšší možná adresa uzlu pro daný prefix,
        např. 192.168.0.255/24
\end{itemize}

\section{Adresy v~IPv6 sítí}\label{ipv6-teorie}

\subsection{Formát adresy}
IPv6 adresa je délky 128 bitů a zapsaných ve formátu {\tt X:X:X:X:X:X:X:X}, kde
{\tt X} je hexadecimální zápis 16 bitového čísla.
Příklad:\\
\begin{verbatim}
FEDC:BA98:7654:3210:fedc:ba98:7654:3210
1080:0000:0000:0000:0008:0800:200C:417a
\end{verbatim}

Preferovaný formát je dále upraven v~RFC 5952\footnote{https://tools.ietf.org/html/rfc5952}, které definuje následující
pravidla:
\begin{itemize}
    \item Nuly na začátku každého 16 bitového čísla je potřeba vynechat (např.
        80 namísto 0080 a 0 namísto 0000).
    \item Více nulových bloků lze nejvýše jednou nahradit znakem {\tt ::}, a
        MUSÍ být nahrazena nejdelší možná posloupnost takových nulových bloků.
        V~případě více shodně dlouhých posloupností, nahrazuje se ta nejvíce
        vlevo.
    \item Znak {\tt ::} nesmí nahrazovat samostatný nulový blok.
    \item Znaky "a", "b", "c", "d", "e", "f" hexadecimální soustavy se vždy
        píšou malými písmeny.
\end{itemize}
Adresy z~předchozího příkladu tedy budou vypadat následovně:\\
\begin{verbatim}
fedc:ba98:7654:3210:fedc:ba98:7654:3210
1080::8:800:200c:417a
\end{verbatim}

Stejně jako v~IPv4 má adresa dvě části: \emph{adresa sítě ({\bf prefix})} a
\emph{adresa uzlu}. Délká prefixu se zapisuje v~desítkovém tvaru za lomítko
(např. {\tt 1080::/60}).

\subsection{Rozdělení adres}
IPv6 adresy je možno rozdělit podle rozsahu. Typicky může jít o~tři
možnosti -- adresy na lince (neprojdou za router), lokální adresy (ULA, nejsou
routovatelné ve veřejné sítí) a veřejné adresy.

\begin{table}[ht!]
    \begin{center}
        \begin{tabular}{l|l|l}
            Význam & Prefix (bitově) & Prefix \\
            \hline\hline
            Veřejné  & {\tt 001} & {\tt 2000::/3} \\
            \hline
            Lokální  & {\tt 1111 110} & {\tt FC00::/7} \\
            \hline
            Linkové  & {\tt 1111 1110 10} & {\tt FE80::/10} \\
            \hline
            Multicast & {\tt 1111 1111} & {\tt FF00::/8} \\
            \hline
        \end{tabular}
    \end{center}
\end{table}

\subsection{Lokální IPv6 (ULA) adresy}\label{ula}
Síťová část ULA adresy se stává ze 4 části:
\begin{description}
    \item [Prefix] {\tt fc00::/7}
    \item [L] bit 1 pokud byl prefix přiřazen lokálně, 0 zatím nebyla
        definována, začátek adresy je proto typicky {\tt fd00::/8}
    \item [Global ID] identifikátor sítě, měl by být unikátní, standard
        popisuje pseudonáhodný algoritmus pro generování. Unikátnost je
        požadována, aby při spojení více lokálních sítí nebylo třeba žádnou
        přečíslovat.
    \item [Inderface ID] Identifikátor rozhraní, existuje několik variant jak
        jej získat. Původní standard (RFC
        3513\footnote{https://tools.ietf.org/html/rfc3513} a RFC
        4291\footnote{https://tools.ietf.org/html/rfc4291}) doporučoval použití
        EUI-64. V~současnosti byla tato metoda nahrazena RFC
        7217\footnote{https://tools.ietf.org/html/rfc7217}. Pro ochranu
        soukromí jsou další specifika automatické generace identifikátorů
        definovány v~RFC 4941\footnote{https://tools.ietf.org/html/rfc4941}.
\end{description}

\begin{table}[ht!]
    \begin{center}
        \begin{tabular}{c|c|c|c|c}
            7 bits & 1 &  40 bits  &  16 bits  & 64 bits \\
            \hline
            1111 110 & L & Global ID & Subnet ID & Interface ID \\
            \hline
        \end{tabular}
    \end{center}
\end{table}

Lokální síť je tedy složená ze tří části:
\begin{itemize}
    \item unikátní prefix délky 48 bitů,
    \item 16 bitový identifikátor podsítě,
    \item adresa uzlu o~délce 64 bitů.
\end{itemize}

\subsection{Manuální konfigurace IP Adres}\label{ip-manual}
Pro dočasnou konfiguraci IP adres na OS Linux slouží příkaz {\tt ip}. Pro
dnešní cvičení budeme potřebovat následující příkazy:
\begin{itemize}
    \item \verb_ip link set [rozhraní] up/down_ pro zapnutí/vypnutí rozhraní
    \item \verb_ip addr add/del [IP adresa]/[prefix] dev [rozhraní]_ pro
        přidání/odstranění adresy z~rozhraní
    \item \verb_ip addr flush dev [rozhraní]_ pro odstranění všech adres
z~rozhraní
    \item \verb_ip route add default via [IP adresa]_ pro nastavení výchozí
        brány
    \item \verb_ip link_, \verb_ip addr_, \verb_ip route_ zobrazí aktuální
        konfiguraci
\end{itemize}

\subsection{Dynamická konfigurace IPv4 - DHCP}\label{dhcp}
Pro dynamickou konfiguraci IP adres je potřeba aby byl na síti přístupný
nakonfigurovaný DHCP server a klienti, kteří si o~IP adresu požádají. Kromě
přiřazení IP adres má DHCP server na starosti i šíření jiných informací
důležitých pro bezproblémovou funkčnost sítě. Jednou takovou informací je IP
adresa doporučeného DNS serveru pro klienty na síti.

Na Vašich počítačích máte nainstalovanou serverovou aplikaci ISC DHCP, která
poskytuje služby DHCP serveru. Aplikace se konfiguruje pomocí souboru
\verb_/etc/dhcp/dhcpd.conf_. Systémová služba se jmenuje
\texttt{isc-dhcp-server.service}.

Příklad obsahu konfiguračního souboru:
\begin{verbatim}
option domain-name-servers [IP adresy DNS serverů];
subnet [IP adresa sítě] netmask [maska] {
    range [první přiřaditelná IP adresa] [poslední přiřaditelná IP adresa];
}
\end{verbatim}
Dodatečné informace o~konfiguračních možnostech najdete v~manuálových stránkách
{\tt man dhcpd.conf} případně {\tt man dhcpd}.

DHCP klient je implementován v~aplikaci \texttt{dhclient}. Může být spuštěn bez
parametrů pro všechna aktivní síťová rozhraní nebo lépe s~jménem konkrétního
síťového rozhraní, které má být konfigurováno:
\begin{verbatim}
dhclient -v [rozhraní]
\end{verbatim}
Argument \texttt{-v} zajistí, že aplikace vypíše detailní informace.
Chybu týkající se {\tt smbd.service} někdy zobrazovanou v~laboratoři můžete ignorovat.

\subsection{Dynamická konfigurace IPv6}\label{dyn-ipv6}
Dynamická konfigurace IPv6 se dělí na {\bf bezstavovou} a {\bf stavovou}
konfiguraci. Na cvičeních se budeme zabývat jen bezstavovou. Bezstavová
konfigurace byla navržena tak, aby stačilo připojit zařízení do sítě a
automaticky si klient vygeneroval nějakou adresu a ihned mohl komunikovat se
světem. K~tomu klient potřebuje znát prefix sítě do které byl připojen.
K~tomuto účelu se používají zprávy označované jako Routing Advertisement (RA).
Tyto a ještě další zprávy jsou součásti procesu Neighbor Discovery (RFC
4861\footnote{https://tools.ietf.org/html/rfc4861}).

Na počítačích v~laboratoři je nainstalovaná systémová služba {\tt radvd.service},
která je schopna vysílat zprávy RA. Aplikaci je možno nakonfigurovat pomocí
souboru {\tt /etc/radvd.conf}.

Příklad konfigurace:
\begin{verbatim}
interface [rozhraní]
{
    AdvSendAdvert on;
    MaxRtrAdvInterval [Max pocet sekund mezi zpravami RA, min 4];
    prefix [prefix]/[delka prefixu]
    {
        AdvOnLink on;
        AdvAutonomous on;
        AdvRouterAddr on;
    };
};
\end{verbatim}
Přehled všech možností konfigurace poskytnou manualové stránky {\tt man radvd}
a {\tt man radvd.conf}.

Pro správné šíření RA zpráv je potřebné povolit směrování IPv6 provozu:
\begin{verbatim}
sysctl net.ipv6.conf.all.forwarding=1
\end{verbatim}

Součásti balíčku radvd je i aplikace {\tt radvdump}, kterou možno použít pro
analýzu -- naslouchá na síťových rozhraních a tiskne na obrazovku obsah
zachycených RA zpráv.

%%%%%%%%%%%%%%%%%%%%%%%%%%%%%%%%%%%%%%
\section{Synchronizace času protokolem NTP (Network Time Protocol)}\label{ntp}
Protokol NTP \cite{rfc1305} umožňuje synchronizovat čas mezi uzly v~síti. Jeho úkolem je přenášet časové informace z primárních časových serverů s přesnými hodinami na ostatní servery a klienty. Protokol NTP při synchronizaci bere v úvahu proměnlivou dobou přenosu paketu, zpoždění paketu a další parametry. NTP organizuje servery hierarchicky do úrovní: primární servery NTP předávají časové informace sekundárním serverům, ty je pak předávají serverům nižší úrovně atd. Úroveň vzdálenosti od přesného zdroje času se vyjadřuje veličinou {\em stratum}. Nejnižší hodnota 0 označuje referenční zdroj přesného času (např. GPS). Stratum 1 označuje primární servery, které jsou synchronizovány právě s~referenčním zdrojem. Stratum 2 jsou servery synchronizovány se servery stratum 1 atd. \cite{rfc5905}.

Implementaci protokolu NTP zajišťuje balík aplikací {\em ntp}. Tento balík se skládá z~několika
aplikací. V~rámci laboratorního cvičení budeme používat aplikace {\tt ntpd} či {\tt ntpq} (viz část \ref{sec:ntpq}). Chrony\footnote{Viz \url{https://chrony.tuxfamily.org/}.} nabízí alternativní
implementaci NTP.

\subsection{Aplikace {\tt ntpd}}
Aplikace {\tt ntpd} (NTP démon) slouží k synchronizaci času pomocí NTP. Může pracovat v režimu klienta i serveru. Aplikace běží na pozadí a neustále synchronizuje čas s~nastavenými servery, případně upravuje lokální čas. Úprava lokálního času se provádí úpravou rychlosti běhu lokálního času. Pokud je lokální čas pozadu, resp. se předbíhá, tak se systémové hodiny {\em zrychlují}, resp. {\em zpomalují}. Tento způsob úpravy času znamená, že pokud je čas odchýlen o~několik minut, bude nějakou dobu trvat, než dojde k synchronizaci. Na druhou stranu se tak zabrání skokové změně a navíc čas se nikdy neposune do minulosti. Pokud je lokální čas odchýlen o~více než 1000 sekund (necelých 17 minut), aplikace to vyhodnotí jako chybu a skončí. Pokud se tak stane, objeví se zpráva v~systémovém logu.

Zda aplikace běží, lze zjistit například příkazem {\tt ntpq -p}. Pokud není démon {\tt ntpd} spuštěn, skončí aplikace hláškou {\em ntpq: read: Connection refused}. 

Službu NTP je možné spustit příkazem:
\begin{verbatim}
systemctl start ntp.service
\end{verbatim}
Aby se předešlo problému v~případě, kdy např. hardwarové hodiny jdou špatně a při vypnutí může dojít k~odchýlení lokálního času o~více než 17 minut, je možné vynutit okamžité nastavení času příkazem
\begin{verbatim}
ntpd -qg  # -g vypne kontrolu odchylky a nastaví čas bez ohledu na velikost odchylky
          # -q ukončí běh démona poté, co byl čas nastaven
\end{verbatim}

\subsubsection{Konfigurace}
Základním konfiguračním souborem je {\tt /etc/ntp.conf}. Tento konfigurační soubor obsahuje mnoho
konfiguračních voleb. Nastavení serverů NTP zajišťuje volba {\tt server}, za níž následuje adresa nebo jméno NTP serveru. Tato volba se může vyskytovat v konfiguračním souboru opakovaně. Pomocí této volby se aplikace {\tt ntpd} přepne do role klienta, který může využít k sychronizaci svého času více serverů NTP. V režimu klient se lokální hodiny synchronizují podle času ze vzdáleného serveru NTP. 

Kromě režimu klient může {\tt ntpd} pracovat v režimu {\em peer}, která umožňuje, aby se vzdálený server synchronizoval podle lokálních hodin. To je užitečné v případě, že je v síti více serverů NTP, které mohou mít lepší zdroj času. Zároveň to řeší problém výpadků serverů. 

V konfiguraci můžeme také zvolit možnosti komunikace typu {\em broadcast} nebo {\em manycastclient}, za nimiž následuje broadcastová či multicastová adresa. Na tyto adresy server NTP posílá periodické zprávy s aktuálním stavem hodin. Pro multicastové přenosy se využívají skupiny definované IANA pro protokol NTP, tj. 224.0.1.1 pro IPv4 a ff05::101 pro IPv6. 

Jinou užitečnou volbou je {\em restrict}, která slouží pro řízení přístupu. Aplikací {\tt ntpd} mohou být zaslány různé požadavky přes síť (například pomocí {\tt ntpq}). Je vhodné povolit některé dotazy jen z~určitého uzlu nebo podsítě. Využití této volby může být také užitečné v~případě, že se pro nastavování času využívá broadcastu nebo multicastu a není žádoucí, aby takto vyslanou informaci klienti akceptovali z~libovolného zdroje. Základní tvar tohoto příkazu je:
\begin{verbatim}
restrict <adresa> [mask <maska>] [<jeden či více příznaků>]
\end{verbatim}

Příznak definuje omezení pro danou adresu/síť:
\begin{itemize}
  \item ignore -- zahazovat všechny pakety
  \item nomodify -- povolí pouze dotazy, požadavky měnící stav serveru jsou zahazovány
  \item noquery -- zakáže dotazy pomocí {\tt ntpq}, synchronizace času není ovlivněna
  \item notrust -- zahazovat neautentizované pakety
\end{itemize}
Další příznaky a jejich popis lze nalézt v~manuálových stránkách {\tt man ntp.conf}.

%%Pokud budou aplikace {\tt ntpq} a {\tt ntpdc} (viz podkapitola \ref{sec:ntpq} používané i pro změnu konfigurace, pak je nezbytné nastavit autentizaci. Protokol nabízí možnost využití symetrické i asymetrické kryptografie. Pro použití symetrické kryptografie jsou k~dispozici tyto volby:
%%\begin{itemize}
%%  \item keys -- tato volba má jeden parametr, který udává název souboru, který obsahuje používané klíče (obvykle {\tt /etc/ntp.keys},
%%  \item trustedkey -- výčet klíčů, kterým se bude důvěřovat,
%%  \item requestkey -- seznam klíčů, které mohou být použity aplikací {\tt ntpdc},
%%  \item controlkey -- seznam klíčů, které mohou být použity aplikací {\tt ntpq}.
%%\end{itemize}
%%Formát souboru {\tt /etc/ntp.keys} má následující tvar:
%%\begin{verbatim}
%%<číslo klíče> <typ> <heslo>
%%\end{verbatim}
%%Typ může nabývat čtyř hodnot. Hodnoty {\tt S} a {\tt N} používají bitový formát a běžně se neužívají. Hodnoty {\tt A} a {\tt M} používají textový řetězce délky 1 až 8 znaků a určují způsob zašifrování při přenosu. Nejčastěji se užívá možnost {\tt M}, která značí použití DES nebo MD5. Příklad souboru pak může vypadat následovně:
%%\begin{verbatim}
%%1 M heslo1
%%2 M secret
%%3 M passwd
%%\end{verbatim}
%%Konfigurace v~{\tt /etc/ntp.conf} potom vypadá:
%%\begin{verbatim}
%%keys /etc/ntp.keys
%%trustedkey 1 2 3
%%requestkey 2
%%controlkey 1 3
%%\end{verbatim}
%%
%%Toto značí, že důvěryhodné jsou všechny tři klíče. Aplikace {\tt ntpdc} se může autentizovat pouzde klíčem 3 a aplikace {\tt ntpq} klíči 1 a 3.   
%%
  
\subsection{Aplikace {\tt ntpq}}\label{sec:ntpq}
%% Tyto dvě aplikace nabízejí v~podstatě podobné možnosti konfigurace serveru NTP. Rozdílů je mezi těmito programy několik. První, který byl již uveden výše, je v~tom, že každá z~těchto aplikací může používat jinou množinu klíčů. Podstatnější rozdíl z~uživatelského hlediska je v~tom, že aplikace {\tt ntpdc} má přesně definovaný výčet příkazů, které lze aplikovat, a proto může měnit jen to, co je v~aplikaci definováno. Naproti tomu {\tt ntpq} disponuje příkazem {\tt :config}, kterému se jako parametry předávají konfigurační volby, které se používají v~souboru {\tt /etc/ntp.conf}. Další rozdíl je ve formátu zpráv, pomocí kterých komunikují se serverem.

Aplikace {\tt ntpq} (NTP query) se používá pro zasílání dotazů serveru NTP {\tt ntpd} z důvodu sledování běhu a výkonnosti, případně zjišťování aktuálního stavu. {\tt ntpq} zasílá zprávy přes síťové rozhraní, i když běží lokálně či vzdáleně. Tato aplikace může měnit konfiguraci {\tt ntpd} za běhu. Možnost změny parametru přes síť může být nežádoucí, proto je vhodné omezit pomocí volby {\em restrict} přístup pouze přes rozhraní {\tt loopback}. 

Příklad výstupu volání {\tt ntpq -p}:

\begin{verbatim}
     remote           refid      st t when poll reach   delay   offset  jitter
==============================================================================
+rhino.cis.vutbr 248.205.243.78   3 u  425 1024  377   10.070   -4.280  10.575
+tik.cesnet.cz   195.113.144.238  2 u  622 1024  377    4.796   -3.464  16.528
*tak.cesnet.cz   .GPS.            1 u  364 1024  377    5.008   -3.625   6.228
\end{verbatim}

Výpis obsahuje následující hodnoty:

\begin{description}

  \item[remote] Klient se synchronizuje vůči třem serverům: rhino.cis.vutbr.cz,
    tik.cesnet.cz a tak.cesnet.cz. Hvězdičkou je označený primární zdroj času,
    plusem sekundární zdroje času.

  \item[refid] Primární zdroj času pro vzdálený server NTP.

  \item[stratum] Počet skoků od vzdáleného serveru k přesnému zdroji času (1
    znamená, že zdroj přesného času je přímo připojen, 16 znamená, že vzdálený
    server je nedosažitelný).

  \item[when] Počet sekund od posledního kontaktu se serverem.

  \item[poll] Počet sekund mezi jednotlivými dotazy protokolem NTP.

  \item[reachability] Úspěšnost posledním 8 pokusů o kontakt vzdáleného serveru
    v osmičkové soustavě. (1 znamená pouze poslední pokus uspěl, 377 znamená
    všech 8 posledních pokusů uspělo).

  \item[delay, offset, jitter] Zpoždění, posun a jitter -- charakteristiky
    vzdáleného zdroje času, podle kterého se volí primární a sekundární zdroje
    času.
\end{description}

%%%%%%%%%%%%%%%%%%%%
\section{Secure Shell}
\label{ssh}

Základní funkcí protokolu Secure Shell (SSH) \cite{rfc4253} je umožnění bezpečného přístupu
 ke vzdálenému počítači přes nezabezpečenou síť. Díky tomu, že je protokol SSH navržen obecně,
 lze pomocí něj zabezpečovat i další služby, jako je např. X Window, přístup ke vzdálenému
 souborovému systému (SFTP, sshfs, scp), tunelování portů TCP apod. Protokol SSH zajišťuje
 šifrování dat, autentizaci, integritu dat a volitelně také kompresi přenášených dat.

V~rámci předmětu ISA se zaměříme pouze na malou část možností, které protokol SSH přináší. V~rámci
 cvičení si vyzkoušíme protokol SSH pro terminálový přístup ke vzdálenému stroji a ukážeme si
 využití přihlašování ke vzdálenému počítači pomocí klíčů. Dále si vyzkoušíme, jak využít SSH
 k~vytvoření jednoduché HTTP proxy.

Jednou z~nejčastěji používaných aplikací pro využití protokolu SSH je sada programů OpenSSH.
 Balíček programů obsahuje kromě klienta {\tt ssh} i serverovou aplikaci {\tt sshd}, program
 pro generování SSH klíčů {\tt ssh-keygen}, agenta pro usnadnění práce s~SSH klíči {\tt ssh-agent}
 a další. My budeme předpokládat, že na počítačích, na které se budeme snažit připojit
 již běží SSH server {\tt sshd}. Konfigurace tohoto programu je nad rámec tohoto manuálu. Zájemci
 mohou nalézt podrobnější informace v~manuálové stránce {\tt sshd\_config(5)}, či v~jiných návodech.

\subsection{Připojení ke~vzdálenému počítači}

Pro připojení se k~počítači pojmenovaném \emph{h01} a otevření příkazového řádku na vzdáleném
 stroji je možné použít příkaz:

\begin{verbatim}
ssh h01
\end{verbatim}

Příkaz ssh má celou řadu parametrů, které jsou detailně popsány v~manuálové stránce {\tt ssh(1)}.
 Z~těch nejčastěji používaných zmíníme alespoň změnu uživatelského jména ({\tt -l}), specifikování
 TCP portu vzdáleného serveru ({\tt -p}), zvýšení výřečnosti programu ({\tt -v}, tento parametr
 lze použít i vícekrát), zapnutí tunelování protokolu X Windows ({\tt -X}, {\tt -Y})
 a přesměrování portů ({\tt -L}). Často zadávané parametry se specifikují v konfiguračním souboru
 \verb|~/.ssh/config|. Popis tohoto souboru obsahuje manuálová stránka {\tt ssh\_config(5)}.

Po připojení ke vzdálenému serveru jsou informace o~použitém spojení dostupné např. v~rámci
 proměnných prostředí. Proměnné prostředí související s~protokolem SSH je možné zobrazit
 příkazem {\tt env | grep SSH}. Následující výpis ukazuje příklad proměnných po připojí
 k~serveru {\tt merlin} protokolem IPv6:

\begin{verbatim}
local $ ssh merlin6.fit.vutbr.cz 
merlin $ env | grep SSH
SSH_CLIENT=2001:67c:1220:80c:e138:4d11:c04c:c675 54514 22
SSH_TTY=/dev/pts/30
SSH_CONNECTION=2001:67c:1220:80c:e138:4d11:c04c:c675 \
  54514 2001:67c:1220:8b0::93e5:b013 22
\end{verbatim}

Z~výpisu vidíme, že připojení bylo realizováno na IPv6 adresu {\tt 2001:67c:1220:8b0::93e5:b013}
 z~počítače s~adresou {\tt 2001:67c:1220:80c:e138:4d11:c04c:c675}. Byl použit zdrojový port č. 54514,
 na straně serveru byl použit standardní protokol 22. Po připojení využíval vzdálený uživatel
 terminál č. 30. Odhlášení ze vzdáleného počítače probíhá standardními prostředky pro ukončení
 shellu, např. příkaz {\tt exit}, nebo vložení konce souboru klávesovou zkratkou {\tt Ctrl-D}.

\subsection{Kopírování souborů mezi počítači}

Pro kopírování souborů protokolem SSH je často využívaná utilita {\tt scp}. Cesta na vzdáleném
 serveru je specifikována v~následujícím formátu: {\tt uživatel@jménoserveru:cesta}.
 Následující příkaz zkopíruje lokální soubor isa na vzdálený počítač h01 do
 složky fit umístěné v~domovské složce uživatele student.

\begin{verbatim}
scp isa student@h01:fit/
\end{verbatim}

\subsection{SSH jako proxy a tunelování portů}
Díky velmi obecné implementaci lze protokolem SSH tunelovat jakýkoli typ provozu. Je tedy možné
 např. zajistit šifrování provozu po určitou část komunikace, konkretně po stanici, ke které
 se připojujeme SSH protokolem, nebo přesměrovat provoz z~některého portu na jiný port. Následujcí
 příkaz nám otevře SSH spojení, které pak můžeme použít jako proxy ve webovém prohlížeči.

\begin{verbatim}
  ssh -D 12345 -N student@sshserver
\end{verbatim}

Po autentizaci stačí pak v~prohlížeči zadat jako nastavení proxy server localhost s~portem 12345
 a~provoz bude směrován nejprve na localhost, poté půjde ssh spojením na sshserver, kde bude převeden
 na běžný HTTP/HTTPS provoz a~poslán dále. 

\subsection{Klíče protokolu SSH}

Klíče protokolu SSH mají několik výhod. Díky nim není nutné posílat heslo pro přístup ke vzdálenému
 stroji přes síť, byť v~šifrované podobě. Délka používaných klíčů znesnadňuje potenciálnímu
 útočníkovi útok hrubou silou, protože síla klíče bývá typicky vyšší než síla běžně používaných
 hesel. Ve spojení s~agentem pro správu klíče může uživatel přistupovat ke vzdálenému serveru bez
 nutnosti opakovaného zadávání hesla. Agent může být nastaven, aby pro použitý klíč nevyžadoval
 heslo po zbytek sezení, či po určitý počet minut.

SSH klíč se obvykle generují utilitou {\tt ssh-keygen}. Po spuštění bez parametrů je vytvořen pár
 klíčů algoritmem RSA o~délce 2048\,B. Uživateli je nabídnuto umístění souboru,
 které může změnit. Dále je uživatel požádán o~zadání passfráze, od které se očekává, že bude
 silnější než heslo. Parametr {\tt -t} nastavuje jiný typ šifrovacího algoritmu,
 parametr {\tt -b} mění délku klíče, parametrem {\tt -f} specifikuje umístění souboru, parametr
 {\tt -C} upravuje popis klíče, další parametry popisuje manuálová stránka {\tt ssh-keygen(1)}.
 Následující příklad vygeneruje klíč o~délce 4096\,B, který nebude chráněn passfrázi:

\begin{verbatim}
ssh-keygen -t rsa -b 4096 -N "" -f ~/.ssh/nopass -C nopass
\end{verbatim}

Po vytvoření klíčů vzniknou dva soubory. Ten který je bez přípony {\tt .pub} je soukromý klíč,
 který by měl zůstat tajný a uživatel, který jej vytvořil by jej neměl dále distribuovat.
 Soubor s~příponou {\tt .pub} je určen pro volnou distribuci.

Podrobný popis použití klíčů během přihlašování najdete např. na \url{https://www.digitalocean.com/community/tutorials/understanding-the-ssh-encryption-and-connection-process}.

\subsection{Konfigurace použití klíčů}

Nejdříve je potřeba distribuovat soubor s~veřejným klíčem na vzdálený počítač. K~tomu slouží
 např. program {\tt scp}. Každý z~uživatelů si může specifikovat sadu klíčů pro přístup k~danému
 stroji v~souboru \verb|~/.ssh/authorized_keys|. Nový klíč do tohoto souboru přidáme např. takto
 (všimněte si, že se klíč přidává na konec souboru pomocí \verb|>>|):

\begin{verbatim}
cat id_rsa.pub >> ~/.ssh/authorized_keys
\end{verbatim}

Nyní se již můžeme připojit ke vzdálenému počítači pomocí vytvořených klíčů. Pokud jsme klíč
 na lokálním počítači umístili do výchozího umístění, je klíč použit automaticky. Pokud jsme
 zvolili jiné umístění, je potřeba program {\tt ssh} informovat o~umístění klíče parametrem
 {\tt -i}, nebo volbou {\tt IdentityFile} v~konfiguračním souboru. Informace o hledaných
 klíčích jsou zobrazeny po použití parametru {\tt -v}.


%%%%%%%%%%%%%%%%%%%%
\section{Transport Layer Security}
\label{tls}

Původní protokoly, na kterých vznikaly původní sítě a počátky Internetu byly
nešifrované. Požadavek na zabezpečení dat se objevil až po vzniku
nejdůležitějších protokolů jako je DNS, Telnet, HTTP, SMTP, FTP apod. Zatímco
protokol SSH popisované v sekci~\ref{ssh} zavedlo nový protokol kompletně
nahrazující protokol Telnet, většina ostatních protokolů založených na TCP
(např. SMTP, HTTP, SIP) funguje v šifrované variantě stejně jako v nešifrované,
jen mezi protokol aplikační vrstvy a TCP přidala mezi vrstvu Secure Sockets
Layer (SSL) a Transport Layer Security (TLS), viz obr.~\ref{fig:tls},
zajišťující šifrování. Výhodou tohoto přístupu je, že přidání podpory šifrované
varianty do existující aplikace je v případě správného návrhu velice snadné a je
potřeba upravit jen funkce zajišťující navazování spojení a zasílání zpráv;
části programu vytvářející a zpracovávající zprávy aplikačního protokolu a
samotný program není potřeba měnit.

\begin{figure}[h!]
  \centering
    \includegraphics[width=0.3\textwidth]{fig/vrstvy-ssl}
  \caption{Vrstvy TCP/IP modelu s využitím SSL/TLS.}
 \label{fig:tls}
\end{figure}

Protokol TLS standardizuje IETF na základě původního protokolu SSL, který je
dnes již zastaralý~\cite{RFC7568}. Poslední verze protokolu TLS je
1.3~\cite{RFC8446}\footnote{Převyprávěná verze RFC 8446 je dostupná na \url{https://www.davidwong.fr/tls13/}}.

TLS vytváří obousměrný šifrovaný kanál mezi dvěma počítači, který zajišťuje:

\begin{itemize}

  \item Autentizaci: vždy se ověřuje identita serveru, identita klienta může být
  také ověřená. Pro ověřování se používají certifikáty X.509, které mohou být
  podepsány certifikační autoritou.

  \item Důvěrnost (utajení): po navázání spojení je obsah přenášené komunikace
  známý pouze komunikujícím stranám. Délka zpráv může být pozorovateli
  komunikace známá, TLS podporuje vycpávkový provoz.

  \item Integritu: Obsah zprávy není možné bez odhalení při přenosu upravit.

\end{itemize}

Samotný protokol TLS se skládá ze dvou podprotokolů:

\begin{itemize}

  \item \emph{Handshake protocol} se používá při navazování spojení a zajišťuje
  autentizaci komunikujících stran, domluvu kryptografických parametrů a
  sdílených klíčů.

  \item \emph{Record protocol} přenáší data mezi komunikujícími stranami.

\end{itemize}

\subsection{Certifikační autority}


Certifikační autorita (CA)\footnote{Text popisující CA je převzatý z~{\tt https://cs.wikipedia.org/wiki/Certifika\%C4\%8Dn\%C3\%AD\_autorita}.} je v asymetrické kryptografii subjekt, který vydává
digitální certifikáty (elektronicky podepsané veřejné šifrovací klíče), čímž
usnadňuje využívání Public Key Infrastructure (PKI) tak, že svojí autoritou
potvrzuje pravdivost údajů, které jsou ve volně dostupném veřejném klíči
uvedeny. Na základě principu přenosu důvěry tak můžeme důvěřovat
údajům uvedeným v digitálním certifikátu za předpokladu, že důvěřujeme samotné
certifikační autoritě.

Na Internetu působí mnoho komerčních certifikačních autorit, které obvykle mají
své veřejné klíče umístěny přímo ve webových prohlížečích a dalších programech,
čímž mohou uživateli zjednodušit rozhodování o míře důvěry webových serverů, ke
kterým se připojuje (ale též digitálně podepsaných e-mailů i jiných dat).
Existují též bezplatné certifikační autority nebo takové, které se řídí zákony
daného státu, vnitřními předpisy organizace a podobně.

Hodnota digitálního certifikátu je úměrná míře důvěry, kterou máme k údajům
v~něm uvedených. Proto je pro certifikační autoritu nejdůležitější důvěra, kterou
vůči svému okolí vzbuzuje (tj. že nevydá digitální certifikát s nepravdivými
údaji). Certifikační autorita proto musí adekvátním způsobem pečovat o svoji
důvěryhodnost, jinak by nebylo možné využít principu přenosu důvěry.

\subsection{Transparentnost certifikátů}

Aby bylo možné ověřit, že CA vydávají certifikáty poctivě, byl experimentálně
zaveden mechanismus transparentnosti certifikátů. Zapojené CA veřejně ohlašují
každý vystavený certifikát (případně precertifikát). Vydané certifikáty se
shromažďují v lozích (cerificate transparency log), kde je monitorovat nově
přidávané certifikáty a auditovat přítomnost certifikátu nalezených při
prohlížení webu.

%%%%%%%%%%%%%%%%%%%%%%%%%%%%%%%%%%%%%%
\section{Domain Name System}
\label{dns}
Cílem DNS je zajistit překlad mezi doménovým jménem a IP adresou. Dříve se jednalo především o~překlad hostname na IP adresu a naopak, tedy záznamy typu {\tt A}, {\tt AAAA}, {\tt PTR} \cite{rfc1035,rfc3596}. Dalším známým typem je {\tt MX}, který deleguje zodpovědnost za příjem e-mailové pošty dané domény. Tento typ se od předchozích tří odlišuje, protože se nejedná o~adresaci hosta, ale služby. Podobný význam mají i záznamy typu {\tt SRV}, které lze v~současnosti využít pro služby SIP a XMPP \cite{rfc2782}.

Mimo tyto typy záznamů existují i další. Některé jsou definovány již v~původním návrhu, jiné přidané později nebo význam původních typů je využit pro další služby (např. {\tt TXT} se používá pro distribuci veřejného klíče podepisování emailu DKIM \cite{rfc4871}).

\subsection{Klient}
Aby klient věděl, na který server se obrátit s~požadavkem na přeložení doménového jména na IP adresu či opačně, je součástí systému tzv. resolver. Tento resolver má několik konfiguračních souborů, z~nichž jmenujme {\tt /etc/hosts}, {\tt /etc/resolv.conf} a {\tt /etc/host.conf}.

První z~těchto souborů se používá pro statický překlad doménového jména na IP adresu. Formát souboru a další popis lze nalézt v~manuálových stránkách {\tt man hosts}.

Dynamický překlad, neboli překlad pomocí DNS serveru, vyžaduje existenci konfiguračního souboru {\tt /etc/resolv.conf}, který obvykle obsahuje jedenkrát volbu {\tt search} a jednu nebo více voleb {\tt nameserver}.
\begin{description}
  \item[search] definuje tzv. searchlist, neboli seznam domén (max. 6), které se budou prohledávat, pokud je požadavek na překlad neúplného doménového jména. Tedy má-li tato volba podobu {\tt search fit.vutbr.cz}, pak při pokusu přeložit jméno {\em merlin} se při neúspěchu pokusí systém také přeložit jméno {\em merlin.fit.vutbr.cz}.
  \item[nameserver] se uvádí právě s~jedním parametrem, který definuje IP adresu DNS serveru, který se systém pokusí kontaktovat. Může se jednat jak o~IPv4 tak o~IPv6 adresu.
\end{description}
Další možnosti, které může soubor obsahovat, lze nalézt v~manuálových stránkách {\tt man resolv.conf}.

V~souboru {\tt /etc/host.conf} lze definovat, v~jakém pořadí se předchozí dvě volby využijí. Standardně se nejprve zkouší statický překlad a poté dynamický. Toto výchozí nastavení umožňuje lokální předefinování překladu z doménového jména na IP adresu.

Pokud se síť na klientovi konfiguruje dynamicky, pak je soubor {\tt /etc/resolv.conf} upravován automaticky. Při statické konfiguraci musí být obsah souboru upraven ručně.

\subsection{Konfigurace serveru}
Konfigurace DNS se stává ze dvou částí -- konfigurace vlastností samotné aplikace a konfigurace obsluhovaných zón. Existuje mnoho různých variant implementací. V~laboratoři se bude pracovat s~implementací od ISC -- BIND.

Konfigurační soubor {\tt\bf named.conf} se na počítačích v~laboratořích nachází ve složce {\tt
/etc/}.

%\begin{verbatim}
%
%acl "local" {
%  <prefix>/<délka prefixu>;
%};
%
%options {
%  listen-on-v6 { <ipv6 adresa lokálního rozhraní>; };
%  allow-query { local; };
%  allow-recursion { local; };
%  allow-transfer { none; };
%  allow-update { none; };
%  directory "/etc/namedb/working";
%  pid-file "/var/run/named/pid";
%}
%
%zone "." IN {
%  type hint;
%  file "/etc/namedb/named.root";
%};
%
%zone "localhost" IN {
%  type master;
%  file "/etc/namedb/master/localhost-forward.db";
%  notify no;
%};
%
%zone "127.in-addr.arpa" IN {
%  type master;
%  file "/etc/namedb/master/localhost-reverse.db";
%  notify no;
%};
%
%zone "0.ip6.arpa" IN {
%  type master;
%  file "/etc/namedb/master/localhost-reverse.db";
%  notify no;
%};
%
%zone "255.in-addr.arpa" IN {
%  type master;
%  file "/etc/namedb/master/empty.db";
%  notify no;
%};
%
%zone "0.in-addr.arpa" IN {
%  type master;
%  file "/etc/namedb/master/empty.db";
%  notify no;
%};
%
%\end{verbatim}

%Zóna typu {\tt hint} odkazuje na seznam tzv. root serverů, tedy takových serverů, které znají informace
%o všech registrovaných top-level doménách a dokáží říci, které servery jsou za tyto top-level domény
%zodpovědné. Zóna typu {\tt master} značí, že server je autoritativním serverem pro danou doménu.
%Mezi povinné záznamy pro takovou doménu patří záznam typu SOA a alespoň jeden NS odkazující na server samotný. Zároveň by měl být obsažen i záznam typu A pro tento server.
Další možnosti konfiguračního souboru naleznete v~manuálových stránkách {\tt man named.conf}.

\subsubsection{Zóna typu hint}
Asociuje se s~nejvyšší doménou v~rámci celé hierarchie. Pokud přijde na server požadavek, prohledá seznam spravovaných zón ostatních typů ({\em master}, {\em slave}, {\em forward}). Pokud záznam nenajde, vybere jeden z~kořenových serverů definovaný právě v~tomto typu zóny.

Soubor, který je vyžadován pro konfiguraci této domény lze získat například programem dig:
\begin{verbatim}
 dig +norec NS . @a.root-servers.net > /var/named/db.root
\end{verbatim}

\subsubsection{Lokálně obsluhovane zóny}
Tyto zóny lze rozdělit do dvou skupin -- {\bf loopback adresy} a {\bf ,,prázdné zóny''}. RFC 6303
\cite{rfc6303} je z~kategorie {\em best practice} a definuje, které zóny by měl být schopen DNS server
obsloužit sám a tím redukovat dotazy přeposílané na další servery. Ve většině případů jsou to zóny pro reverzní záznamy privátních IP adres.

%Výchozí soubor aplikace {\bf named} ve FreeBSD respektuje právě RFC 6303 a navíc doplňuje ještě další
%zóny. Například nealokované rozsahy IPv6 adres. Pro práci v~laboratoři nahraďte původní soubor {\tt
%/etc/namedb/named.conf} souborem {\tt /root/isa1/named.conf}, který je zkrácen a obsahuje pouze nezbytné definice.

\subsubsection{Vlastní zóna}\label{sec:vlastni_zona} 
Aby server začal překládat vlastní zónu, je třeba přidat definici zóny do souboru {\tt /etc/named.conf} a vytvořit zónový soubor. Pokud se přidává zónový soubor pro doménu, je vhodné přidat i zónu pro reverzní záznamy.

\begin{verbatim}
zone "moje.domena.cz" {
  type master;
  file "/var/named/moje.domena.cz.zone";
};

zone "0.168.192.in-addr.arpa" {
  type master;
  file "/var/named/192.168.0.rev";
};
\end{verbatim}

Pro každou z~těchto zón se tvoří samostatný zónový soubor, jehož základní tvar má podobu:

%$ORIGIN <nazev domény>
\begin{verbatim}
$TTL 5m

@ IN  SOA <fqdn autoritativniho serveru>. <email správce>. (
  1   ; seriové číslo
  10h ; obnovovací interval pro sekundární servery
  10m ; prodleva, po které se sekundární server pokusí znova kontaktovat
      ; autoritativní server, pokud se předchozi spojení nezdařilo
  1w  ; jestliže se sekundárnímu serveru nepodaří kontaktovat autoritativní
      ; server během této doby, přestane vracet záznamy pro tuto doménu
  1h  ; původně výchozí TTL (nahrazeno $TTL), nově určuje
      ; NEGATIVNÍ TTL -- pro chybové odpovědi (např. NXDOMAIN)
  )

@ IN  NS <fqdn autoritativního serveru>
@ IN  NS <fqdn sekundárního serveru>

následuje seznam dalších záznamů
\end{verbatim}
Typ záznamu {\tt\bf SOA} definuje parametry zóny -- jméno autoritativního serveru a email správce domény. Všimněte si, že znak `{\tt @}' má zvláštní význam (zastupuje název domény). Z~tohoto důvodu se v~emailu správce nahrazuje znak `{\tt @}' za znak `{\tt .}' (tečka).

Záznam typu {\tt\bf NS} by měl být vždy obsažen alespoň jednou a to právě pro autoritativní server. Počet sekundárních serverů je libovolný. Tyto záznamy by také měly byt obsaženy v~nadřazené doméně, aby bylo možné nalézt server zodpovědný za danou poddoménu. Tedy kořenová doména obsahuje {\tt NS} záznamy pro domény nejvyšší úrovně {\em .cz, .com, .eu, \dots} a ty zase pro jednotlivé poddomény {\em google.com, seznam.cz, \dots}.

V~zónovém souboru se rozlišuje mezi relativním a absolutním jménem. O~jaký typ jména jde, se rozlišuje podle zakončení. Končí-li tečkou, je to jméno absolutní (plně kvalifikované). V~opačném případě jde o~jméno relativní a za to se vždy doplní obsah definován v~{\tt \$ORIGIN}. Jestliže není definován, použije se název, který je definován u~názvu zóny v~souboru {\tt named.conf}.

Další záznamy mají obecný tvar:
\begin{verbatim}
jméno  ttl  třída  typ  <na typu závislá data>
\end{verbatim}
Jméno běžně bývá ve tvaru relativním. Pokud není hodnota TTL jiná než je výchozí pro celou zónu, pak se může vynechat. Třída záznamu se běžně používá pouze {\tt IN} (Internet). Kromě typu SOA a NS, se v~laboratorním cvičení použijí záznamy typu {\tt A}, {\tt PTR} a volitelně {\tt AAAA}. Typ {\tt AAAA} se chová stejně jako záznam typu {\tt A}, který se používá pro IPv4. Typově závislými daty je v~tomto případě adresa IPv6, resp. IPv4. Příklad:
\begin{verbatim}
host  IN A  192.168.0.1
\end{verbatim}
\begin{verbatim}
host  IN AAAA  fd12:1234::1
\end{verbatim}
Typ {\tt PTR} nerozlišuje, zda se jedná o~IPv6 nebo IPv4 adresu. Pro tento případ se typově závislá data chovají stejně jako jméno na začátku záznamu. V~tomto případě se ale vždy zapisuje v~plném tvaru. Další zvláštností PTR záznamu, který je vidět i v~názvu zóny, je, že se hodnoty zapisují v~obráceném pořadí a každá hodnota je oddělená tečkou. Obrácené pořadí je kvůli vyhodnocování záznamů od konce a tečka mezi každým číslem dovoluje snadnější delegaci podsítě na jiný DNS server. Příklad:
\begin{verbatim}
1		IN	PTR	host.moje.domena.cz.
\end{verbatim}

Pro IPv6 můžete předefinovat proměnnou {\tt \$ORIGIN} např. následovně:
\begin{verbatim}
$ORIGIN 0.0.0.0.0.0.0.0.4.3.2.1.2.1.d.f.ip6.arpa.
1.0.0.0.0.0.0.0.0.0.0.0.0.0.0.0   IN PTR host.moje.domena.cz.
\end{verbatim}
Pozor však na to, aby expanzí nedocházelo k vytváření neplatných záznamů jako jsou
\begin{verbatim}
$ORIGIN 0.0.0.0.0.0.0.0.4.3.2.1.2.1.d.f.ip6.arpa.
2.0.0.0.0.0.0.0.0.0.0.0.0.0.0.0   IN PTR host
1.0.0.0.0.0.0.0.0.0.0.0.0.0.0.0   IN PTR host.moje.domena.cz
\end{verbatim}
protože po expanzi se z nich stanou následující záznamy:
{\small
\begin{verbatim}
2.0.0.0.0.0.0.0.0.0.0.0.0.0.0.0.0.0.0.0.0.0.0.0.4.3.2.1.2.1.d.f.ip6.arpa. IN PTR
                                              host.0.0.0.0.0.0.0.0.4.3.2.1.2.1.d.f.ip6.arpa.
1.0.0.0.0.0.0.0.0.0.0.0.0.0.0.0.0.0.0.0.0.0.0.0.4.3.2.1.2.1.d.f.ip6.arpa. IN PTR
                               host.moje.domena.cz.0.0.0.0.0.0.0.0.4.3.2.1.2.1.d.f.ip6.arpa.
\end{verbatim}}

Příklad PTR záznamu pro IPv4 je:
\begin{verbatim}
19.176.229.147.in-addr.arpa. IN	PTR	merlin.fit.vutbr.cz.
\end{verbatim}

Po každé změně zónového souboru by se vždy mělo upravit i sériové číslo zóny a to tak, aby bylo vždy větší než předchozí hodnota. Také je nutné aplikaci načíst změněnou konfiguraci.
\begin{verbatim}
systemctl restart named.service
\end{verbatim}

\subsection{DNS Over TLS}\label{DoT}
DNS Over TLS (DOT) přináší šifrování a autentizaci do DNS komunikace. Tento protokol je definován v RFC 7858\cite{RFC7858}. Protokol má přidělen standardizovaný port 853 jako defaultní. TLS protokol zde zaobaluje DNS požadavky a odpovědi jako přenášený payload viz obrázek \ref{fig:dot-encapsulation}.

\begin{figure}[h]
    \centering
    \includegraphics[scale=1]{fig/dot.pdf}
    \caption{DNS Over TLS zapouzdření.}
    \label{fig:dot-encapsulation}
\end{figure}

Komunikace probíhá pomocí klasického klient-server schématu a transportním protokolem je TCP. Nejprve proběhne TCP handshake, následuje TLS handshake tak jak je definován v RFC 5246\cite{RFC5246}. Jakmile je spojení ustaveno dojde k přenosu dat.

Pro minimalizaci latence klienti nemusí čekat na odpovědi serveru, ale můžou použít pipelining pro odesílání více dotazů v rámci jednoho TLS spojení. Přenášená data uvnitř TLS jsou stejného formátu jako v případě standardního DNS přenášeného přes TCP transportní protokol.\cite{RFC7858}

Resolvíng pomocí DNS over TLS lze na Linuxových systémech nastavit například v podobě spuštění Unbound DNS Caching serveru. Kde je nutné upravit konfigurační soubor \texttt{/etc/unbound/unbound.conf} tak, aby obsahoval následující:

\begin{verbatim}
forward-zone:
    name: "."
    forward-ssl-upstream: yes
    # Cloudflare DNS
    forward-addr: 1.1.1.1@853
    forward-addr: 1.0.0.1@853
\end{verbatim}

Následně službu zapnout a nastavit systémový resolver na adresu \texttt{127.0.0.1}.

\subsection{DNS Over HTTPS}\label{DoH}
DNS Over HTTPS (DOH) je nový standard pro bezpečný a šifrovanz přenos DNS provozu tunelovaný v protokolu HTTPS. Tento protokol nemá vlastní definovaný port a je přenášen jako ostatní HTTPS provoz pod portem 443. DOH je definován standardem vydaným roku 2018 skrývajícím se pod dokumentem označeným jako RFC 8484 \cite{RFC8484}.

\begin{figure}[h]
    \centering
    \includegraphics[scale=0.75]{fig/doh-encapsulation.pdf}
    \caption{DNS over HTTPS zapouzdření.}
    \label{fig:doh-encapsulation}
\end{figure}

DNS zprávy jsou zde přenášeny standardními metodami HTTP GET a POST. V případě GET je DNS request zakódován ve formátu URL64Base jako parametr URI. Výsledná odpověď je obdržena v těle HTTP odpovědi ve standardním DNS wire formátu. V případě metody POST je již sám dotaz zakódován v DNS wire formátu a přenášen v těle POST dotazu.


%%%%%%%%%%%%%%%%%%%%%%%%%%%%%%%%%%%%%%
\section{Přenos systémových události službou Syslog}\label{syslog}
Služba Syslog  slouží pro generování a přenos systémových událostí síťových zařízení a koncových stanic na centrální úložiště (server). Systémové událostí jsou generovány operačním systémem či běžícími aplikace. Tyto události se lokálně ukládají do logovacích souborů. Služba Syslog posílá vygenerované události na server. 

Architekturu služby Syslog tvoří protokol Syslog \cite{rfc5424}, který popisuje formát přenosu zpráv, dále Syslog klient ({\tt originator}), které vytvořené události posílá na server, a Syslog server ({\tt collector}), který přijímá poslané události protokolem Syslog od klientů a centrálně je zpracovává. 

Každá vygenerovaná událost se ukládá jako záznam do logovacího souboru buď operačního systému (např. \verb|/var/log/messages|) nebo do logovacího souboru konkrétní aplikace (např. \verb|/var/log/maillog| pro poštovní server, \verb|/var/log/ipfw| pro firewall, \verb|/var/log/apache/access.log| pro HTTP server).  

Logovací záznam události zpravidla obsahuje časovou značku, typ zařízení, které událost vygenerovalo (facility), závažnost události (severity) a slovní popis události. V~některých logovacích souborech se může objevit i doménové jméno či IP adresa stanice, která událost způsobila. Záznamy mohou být uloženy lokálně v souborovém systému (na unixových systémech v adresáři \verb|/var/log|) nebo je můžeme přeposílat protokolem Syslog na vzdálený server. 

Typ zařízení, které událost vygenerovalo, je definováno identifikátorem v záznamu události. Seznam zařízení, které definuje protokol Syslog, je v tabulce \ref{tab:facility}. 
\begin{table}[h]
  \centering
  \small
  \begin{tabular}{|c|l|c|l|}
    \hline
    \bf Code & \bf Facility & \bf Code & \bf Facility\\
    \hline
     0  &  kernel messages                             &  12 &  NTP subsystem           \\
     1  &  user-level messages                         &  13 &   log audit              \\
     2  &  mail system                                 &  14 &   log alert              \\
     3  &  system daemons                              &  15 &   clock daemon (note 2)  \\
     4  &  security/authorization messages             &  16 &   local use 0  (local0)  \\
     5  &  messages generated internally by syslogd    &  17 &   local use 1  (local1)  \\
     6  &  line printer subsystem                      &  18 &   local use 2  (local2)  \\
     7  &  network news subsystem                      &  19 &   local use 3  (local3)  \\
     8  &  UUCP subsystem                              &  20 &   local use 4  (local4)  \\
     9  &  clock daemon                                &  21 &   local use 5  (local5)  \\
     10 &  security/authorization messages             &  22 &   local use 6  (local6)  \\
     11 &  FTP daemon                                  &  23 &   local use 7  (local7)  \\
     \hline
  \end{tabular}
  \caption{Typy zařízení Syslog dle RFC 5424 \cite{rfc5424}}\label{tab:facility}
\end{table}
Protože tento seznam vychází z původního  návrhu protokolu BSD Syslog z roku 2001, obsahuje zařízení, která se dnes už příliš nepoužívají. Z tohoto důvodu se pro aplikace a další typy zdrojů událostí používá lokální označení s  kódy 16-23. 

Pro hodnocení závažnosti vzniklé události definuje Syslog osm úrovní od nejzávažnější typu Emergency (0) po ladící výpisy typu Debug (7). Pro správu sítí jsou důležité zejména události od úrovně Warning (4). Přehled definovaných úrovní závažnosti je v tabulce \ref{tab:severity}. 
\begin{table}[h]
  \centering
  \small
  \begin{tabular}{|c|l|l|}
    \hline
    \bf Code & \bf Severity & \bf Description \\
    \hline
    0 & EMERGENCY & system is unusable  \\
    1 & ALERT & action must be taken immediately\\
    2 & CRITICAL & critical conditions \\
    3 & ERROR & error conditions\\
    4 & WARNING & warning conditions\\
    5 & NOTICE & normal but significant condition\\
    6 & INFORMATIONAL & informational messages\\
    7 & DEBUG & debug-level messages\\
    \hline
  \end{tabular}
  \caption{Úrovně závažnosti událostí Syslog dle RFC 5424 \cite{rfc5424}}\label{tab:severity}
\end{table}


%%%%%%%%%%%%%%%%%%%%%%%%%%%%%%%%%%%%%%
\clearpage
\section{Monitorování sítě SNMP}\label{snmp}

Služba SNMP (Simple Network Monitoring Protocol) slouží k aktivnímu monitorování zařízení na síti. Protokol SNMP \cite{rfc3412} slouží k přenosu monitorovacích informací mezi řídící stanicí NMS (Network Management Station) a monitorovaným zařízením (SNMP Agent). Protokol slouží nejen ke čtení provozních parametrů (hodnot objektů MIB) připojených stanic (využití paměti, počty přenesených paketů, stav síťových rozhraní, zatížení CPU apod.), ale dá se použít i k  nastavení některých těchto parametrů. Protokol pracuje nad UDP na portech 161 (vyzývání) a 162 (asynchronní zprávy, traps).

Řídicí stanice NMS posílá zprávy SNMP na jednotlivá síťová zařízení s cílem získat hodnoty sledovaných objektů systému. Standardní způsob komunikace pracuje na principu vyzývání (polling), kdy řídicí stanice NMS periodicky posílá zprávy typu dotaz na sledovaná zařízení. Požadované hodnoty (objekty MIB) jsou identifikovány pomocí tzv. OID (Object Identifier), které jsou standardizovány v databázi MIB (Management Information Base).

Příkladem objektu je {\tt sysDescr} (System Description) s OID {\tt 1.3.6.1.2.1.1.6.0}. Standardizované objekty jsou uloženy v databázi MIB-2, která je definována standardem RFC 1213 \cite{rfc1213}.

Pro vyhledávání objektů SNMP lze využít následující portály:
\begin{itemize}
  \item Global OID Reference Database: \url{http://oidref.com/}
  \item Object Identifier (OID) Repository: \url{http://www.oid-info.com/}
\end{itemize}

\subsection{Nástroj pro monitorování zařízení SNMP}
Pro monitorování zařízení SNMP můžeme využít nástroj {\tt snmpwalk}, viz {\tt man snmpwalk}, který slouží k získání objektů z daného zařízení protokolem SNMP. Program {\tt snmpwalk} vrátí buď požadovaný objekt MIB zadaný identifikátorem OID nebo množinu objektů MIB, které patří do zadaného podstromu MIB.

Formát a příklad příkazu {\tt snmpwalk}:
{\small
\begin{verbatim}
snmpwalk -v <protocol> -c <community string> <host> <object>

snmpwalk -v2c -c public isa.fit.vutbr.cz sysDescr         ; vyhledá objekt System Description
snmpwalk -v2c -c public isa.fit.vutbr.cz 1.3.6.1.2.1.1.1  ; ekvivalentní zápis pomocí OID

snmpwalk -v2c -c public isa.fit.vutbr.cz system           ; vyhledá objekty podstromu System
snmpwalk -v2c -c public isa.fit.vutbr.cz 1.3.6.1.2.1.1    ; ekvivalentní zápis pomocí OID
\end{verbatim}
}

\subsection{Konfigurace agentů SNMP}
Zařízení v síti SNMP jsou rozdělena do logických celků, komunit, identifikovaných textovým  řetězcem. Komunita slouží k autentizaci zařízení v síti a definuje přístup k informacím MIB.

V konfiguraci monitorovaného zařízení (agenta SNMP) můžeme definovat, kdo může přistupovat k jakým objektů SNMP. Můžeme definovat způsob přístupu k objektům SNMP (pro čtení: {\tt ro},  pro čtení a zápis: {\tt rw}), můžeme omezit přístup na určité zdrojové adresy, či můžeme definovat část podstromu MIB, do kterého má daný uživatel přístup.  
Příklad nastavení: 
\begin{verbatim}
rocommunity public 10.0.0.0/24
rocommunity public 192.168.0.0/24 1.3.6.1.2.1.1
rwcommunity public localhost
\end{verbatim}
Výše uvedená konfigurace umožňuje zařízením v síti 10.0.0.0/24 přístup k objektům MIB komunity {\tt public} pouze
 pro čtení. Stanice v síti 192.168.0.0/24 mají přístup pro čtení omezen pouze na základní systémové informace uložené v podstromu MIB-2 {\tt system} (OID 1.3.6.1.2.1.1). 
 (podstrom system). Pro lokálního uživatele budou objekty MIB přístupné pro čtení i pro zápis.

 %%%%%%%%%%%%%%%%%%%%%%%%%%%%%%%%%%%%%%
 \newpage
\section{Monitorování toků NetFlow}\label{netflow}
Služba NetFlow byla navržena pro monitorování toků v síti. Její architekturu tvoří sonda NetFlow (exportér), která monitoruje síťový provoz, analyzuje jednotlivé pakety na úrovni L3 a L4, a vytváří tzv. záznamy (flow records) o probíhajícím provozu. Záznamy jsou uložený v lokální paměti sondy, který je po skončení toku přepošle protokolem NetFlow na tzv. kolektor. Kolektor ukládá záznamy z různých sond do centrálního úložiště, kde je můžeme vyhledávat, analyzovat, identifikovat potenciálně nebezpečný provoz apod.

Příklad záznamu o toku:
{\footnotesize
\begin{verbatim}
Date flow start   Duration Proto Src IP Addr:Port      Dst IP Addr:Port   Flags Tos Packets Bytes Flows
2005-08-30 06:53:53 63.545 TC  113.138.32.152:15225 ->  222.33.70.124:3575.AP.SF 0    62     3512   1
2005-08-30 06:53:53 63.545 TC  222.33.70.124:3575   ->  113.138.32.152:25 .AP.SF 0    58     3300   1
\end{verbatim}
}
Tento záznam obsahuje dva toky TCP, z nichž první tvoří 62 paketů o velikosti 3.512 B, které byly poslány z adresy 113.138.32.152 a portu 15225 an adresu 222.33.70.124, port 3575. Délka přenosu trvala 63,545 sekund. 


\subsection{Program {\tt nfdump}}
Program {\tt nfdump} slouží jako kolektor systému NetFlow, viz {\tt man nfdump}. Program také umožňuje zobrazovat uložené záznamy, vyhledávat v nich, provádět agregaci, filtrování a další. 

\begin{verbatim}
  Format:   nfdump [options] [filter] 
\end{verbatim}

Příklady použití:
\begin{verbatim}
nfdump -r nfcapd.202004221040  # vypíše všechny záznamy o tocích z daného souboru

nfdump -r nfcapd.202004221040 'host 147.229.5.2'  # vypíše všechny toky dané stanice
 
nfdump -r nfcapd1 -c 100 'proto tcp and (src ip 172.16.17.18 or dst ip 172.16.17.19)'
         # vybere ze souboru 100 záznamů, které vyhovují zadanému filtru

nfdump -r nfcapd1 -s record/bytes -n 10
         # vygeneruje statistiku Top 10 záznamů s největším počtem přenesených bytů

nfdump -r nfcapd1 -s srcip/packets -n 10
         # vygeneruje statistiku Top 10 IP adres s největším počtem odeslaných paketů

nfdump -r nfcapd1 -A srcip,dstport
         # vypíše agregované toky podle zdrojové IP adresy a cílového portu

nfdump -r nfcapd1 -A srcip 'host 147.229.5.2'
         # vypíše agregované toky podle zdrojové IP adresy pro daný filtr

nfdump -r nfcapd1 -t 2020/04/22.12:19:00-2020/04/22.12:20:00 'port 53'
         # vyhledá všechny záznamy s komunikací DNS v daném časovém okně 
\end{verbatim}

%% 
%% Formát záznamu NetFlow umožňuje ukládání informací o IP provozu na síťovém zařízení. Formát byl
%%  vytvořen firmou Cisco jako funkcionalita jejich zařízení. Ve verzi 9 pak byl standartizován
%%  jako RFC 3954 \cite{rfc3954}.
%% 
%% NetFlow záznamy jsou využívány k~monitorování a analýze provozu v~síti. Data obsahují záznamy
%%  o IP adresách, portech, počty přenesených paketů či bajtů apod. Data jsou uchovávána pro
%%  tzv. toky neboli flow. Toky jsou identifikovány pěticí zdrojová a cílová IP adresa,
%%  zdrojový a cílový port a protokol transportní vrstvy. Tok je chápán jako jednosměrná komunikace.
%% 
%% Infrastruktura pro sběr a uchovánvání NetFlow dat sestává z~kolektoru a sond. Kolektorem je míněna
%%  stanice, která sbírá data ze sítě a provádí jejich uchování. Sondou pak může být jakékoli zařízení
%%  ať už směrovač, přepínač nebo zařízení přímo vyhrazené pro tuto činnost. Sběr dat může probíhat
%%  pro všechny toky, nebo může být prováděno vzorkování či agregace. 
%% 
%%%%%%%%%%%%%%%%%%%%%%%%%%%%%%%%%%%%%%

%%%%%%%%%%%%%%%%%%%%%%%%%%%%%%%%%%%%%%
\section{Signalizační protokol SIP}\label{sip}
Session Initiation Protocol (SIP, RFC 3261 \cite{rfc3261}) je signalizační protokol pro IP telefonie, který provádí následující operace:
\begin{itemize}
    \item registrace uživatel VoIP na serveru SIP
    \item navázání hovoru, výměna parametrů spojení
    \item směrování hovorů
    \item ukončení či zrušení hovoru
\end{itemize}

Protokol SIP neprovádí správu relací po jejich navázání, nezajišťuje kvalitu přenosu, neslouží k přenosu hlasových dat.

Základními prvky IP telefonie jsou {\em server UAS} (User Agent Server), který provádí registraci klientů, navazování a směrování hovorů, lokalizaci klienta účtování apod., a {\em uživatelský agent UAC} (User Agent Client), který zprostředkovává uživateli volání (vytáčení spojení, příjem telefonátu). 

Služba SIP používá pro adresování speciální URI (Uniform Resource Identifier) ve tvaru {\tt sip:user @domain}, kde {\tt user} je uživatelské jméno, které přidělí poskytovatel uživateli, a {\tt domain} je VoIP doména poskytovatele. Adresa uživatele je součástí metody {\tt INVITE}, která slouží k vytvoření hovoru, nebo se využívá při registraci telefonu (metoda {\tt REGISTER}). 

Směrovací informace se během navazování spojení ukládají v hlavičce SIP paketu {\tt Via:}, {\tt Route} či {\tt Record-Route}. 

\subsection{Příklad signalizace SIP}
{\footnotesize
\begin{verbatim}
INVITE sip:541141118@cesnet.cz SIP/2.0    # navázání spojení s uživatelem 541141118 v doméně cesnet.cz
Call-ID: D40CA785-2EEE-4801-9B04-349632F56CDC@147.229.14.146   # jednoznačný identifikátor hovoru
CSeq: 2 INVITE                                                 # sekvenční číslo zprávy, název metody
From: "Petr Matousek"<sip:matousp@cesnet.cz>;tag=2007034328740 # identifikace volajícího (SIP URI)
To: <sip:541151118@cesnet.cz>                                  # identifikace volaného (SIP URI)
Contact: <sip:matousp@147.229.14.146:59183;ob>                 # identifikace volaného (Device URI)

SIP/2.0 100 Trying -- your call is important to us             # odpověď na metodu INVITE
Call-ID: D40CA785-2EEE-4801-9B04-349632F56CDC@147.229.14.146   # jednoznačný identifikátor hovoru
CSeq: 2 INVITE                                                 # odpověď na metodu INVITE   
From: "Petr Matousek"<sip:matousp@cesnet.cz>;tag=2007034328740 # zopakování hlaviček z metody INVITE
To: <sip:541151118@cesnet.cz>
\end{verbatim}
}

\subsection{Základní metody protokolu SIP}
\begin{itemize}
    \item REGISTER -- žádost o registraci
    \item INVITE, ACK -- vytváření spojení
    \item CANCEL -- zrušení vytvářeného spojení
    \item BYE -- ukončení spojení
    \item OPTIONS -- zjišťování možností přenosu
    \item PRACK -- provizorní ACK
\end{itemize}

\subsection{Další metody protokolu SIP}
\begin{itemize}
    \item INFO (RFC 6086) – přenost aplikačních informací
    \item MESSAGE (RFC 3428) – textové zprávy IM
    \item SUBSCRIBE, NOTIFY (RFC 3265) – zasílání upozornění na události
    \item PUBLISH (RFC 3903) – zveřejnění stavu událostí (prezence)
\end{itemize}

\subsection{Položky hlavičky SIP}
\begin{itemize}
    \item {\bf Via:} obsahuje směrovací informace. Každý server, přes který zpráva {\tt INVITE} prochází,
      přidá do hlavičky jeden řádek {\tt Via:} s IP adresou a portem, na kterém zpracoval paket. Zaznamenaná posloupnost uzlů, přes které zpráva {\tt INVITE} procházela, se pak využije pro směrování odpovědi. 
    \begin{verbatim}
    via: SIP/2.0/UDP 10.11.12.51:5062;rport;branch=z9hG4bk2079031997
    \end{verbatim}
    \item {\bf From:} obsahuje odesílatele SIP zprávy ve formátu SIP URI. Nepovinně může obsahovat i zobrazitelné jméno (display-name), které uživatel vložil do konfigurace klienta. 
    \begin{verbatim}
    From: "user08" <sip:user08@iss.cz>;tag=283687066
    \end{verbatim}
  \item {\bf To:} obsahuje adresu příjemce zprávy. 
    \begin{verbatim}
    To: <sip:1006@iss.cz>
    \end{verbatim}
    \item {\bf Allow} obsahuje seznam metod SIP, které podporuje daný klient.
    \begin{verbatim}
    Allow: INVITE, INFO, PRACK, ACK, BYE, CANCEL, OPTIONS, NOTIFY, REGISTER
    \end{verbatim}
\end{itemize}

\subsection{Protokol SDP (Session Description Protocol)}
Protokol SDP \cite{rfc4566} se používá pro zaslání informací potřebných k vytvoření kanálu pro přenos hlasu (případně obrazu). Jde opět o signalizační protokol, který domlouvá parametry spojení, jako je typ média (video, audio, atd.), transportní protokol (RTP/UDP/IP, H.320, atd.), typ kodeku, IP adresa a port pro příjem hlasu (obrazu).
Zprávy SDP se přenáší ve zprávách typu {\tt INVITE} a {\tt OK}.
{\footnotesize
\begin{verbatim}
v=0                                             # verze protokolu SDP
o=- 3434794301 3434794301 IN IP4 147.229.14.146 # původce spojení (username, IP adresa)
s=Sjphone                                       # jméno spojení: jméno softwaru
c=IN IP4 147.229.14.146                         # spojení: IP adresa pro příjem dat
t=0 0 # time (start + end)
a=direction:passive                               
m=audio 49152 RTP/AVP 3 97 98 8 0_101           # audio stream: číslo portu, nabízené kodeky
a=rtpmap:3 GSM/8000
a=rtpmap:97 iLBC/8000
a=rtpmap:98 iLBC/8000
a=fmtp:98 mode=20
a=rtpmap:8 PCMA/8000
a=rtpmap:0 PCMU/8000
a=rtpmap:101 telephone-event/8000
a=fmtp:101 0-11,16
\end{verbatim} 
}


\newpage

%% BIBLIOGRAPHY
\bibliography{../rfc,manual}
\bibliographystyle{plain}

\newpage
\thispagestyle{empty}

\end{document}
%% END OF FILE 
